% ===================================================================
%  QUADRO N.º 11 — A Cidade e os seus monumentos. — O Incêndio.
%
%  Ceci n'est PAS une transcription : c'est une traduction, faite en
%  2026, en portugais d'aujourd'hui. Voir texto/pt/10-quadro-01.tex
%  pour la consigne, qui vaut pour les seize fichiers.
% ===================================================================

%%K t11-apar-1 apar
\VUsaut{2.0mm}
\VUcentre{13.2pt}{90}{TERCEIRA SÉRIE}
\VUsaut{3.0mm}
\VUfilet{16mm}
\VUsaut{5.6mm}
\VUtitre{15.0pt}{60}{QUADRO N.º 11}
\VUsaut{1.4mm}
\VUpk{11.4pt}{40}{A Cidade e os seus monumentos. --- O Incêndio.}
\VUsaut{2.2mm}

%%K t11-tit-1 sub
\VUpk{10.2pt}{40}{Os arredores.}
\VUsaut{3.0mm}

%%K t11-c1-01-1 p
1. --- Moro numa grande cidade situada a cem quilómetros do oceano e que,
no entanto, é um \VUgras{porto} \textsuperscript{(1)} de primeira ordem. O rio que
a banha tem quatrocentos metros de largura e forma uma enseada muito
bela.

%%K t11-c1-02-1 p
2. --- Toda a parte da cidade que está junto aos \VUgras{cais}
\textsuperscript{(2)} tem um aspeto completamente marítimo e industrial e forma
um \VUgras{arrabalde} \textsuperscript{(3)} habitado só por marinheiros e
operários. Estes trabalham nas \VUgras{fábricas} \textsuperscript{(4)} e na
\VUgras{fábrica de gás} \textsuperscript{(5)}, cuja alta \VUgras{chaminé}
\textsuperscript{(6)} e cujo \VUgras{gasómetro} se veem de muito longe.

%%K t11-c1-03-1 p
3. --- Na Idade Média a cidade estava fortificada, e ainda se encontram
alguns restos das suas muralhas, em particular a \VUgras{porta}
\textsuperscript{(8)} do velho \VUgras{castelo forte} \textsuperscript{(7)}, ladeado
pelas suas duas \VUgras{torres} \textsuperscript{(9)}, que hoje servem de
\VUgras{cadeia}.

%%K t11-c1-04-1 p
4. --- Em todo este \VUgras{bairro}, de aspeto muito antigo, um só
edifício é relativamente moderno: a \VUgras{alfândega} \textsuperscript{(10)}. Está
entre o cais e a pequena \VUgras{rua} \textsuperscript{(11)} que leva à velha
\VUgras{casa da câmara} \textsuperscript{(12)}. Este último monumento reconhece-se
facilmente pelo gigantesco \VUgras{campanário} \textsuperscript{(13)} que domina
toda a cidade velha.

%%K t11-c1-05-1 p
5. --- Do outro lado da rua \VUgras{ergue-se} um edifício construído no
mesmo estilo que a alfândega: é a \VUgras{bolsa} \textsuperscript{(14)}. Uma das
suas fachadas dá para uma pracinha na qual está o \VUgras{governo civil}
\textsuperscript{(15)}. A nossa cidade, com efeito, sem ser uma capital, é uma
importante cabeça de província.

%%K t11-tit-2 sub
\VUsaut{5.0mm}
\VUpk{10.2pt}{40}{A Parte alta da Cidade.}
\VUsaut{3.4mm}

%%K t11-c1-06-1 p
6. --- A guarnição, bastante numerosa, aloja-se em vastos
\VUgras{quartéis} \textsuperscript{(16)} muito salubres; estendem-se até à grade do
\VUgras{parque municipal} \textsuperscript{(17)}. A \VUgras{avenida}
\textsuperscript{(19)} que leva a esse parque passa por trás da \VUgras{caixa
económica} \textsuperscript{(18)} e desemboca na praça da \VUgras{catedral}
\textsuperscript{(20)}.

%%K t11-c1-07-1 p
7. --- Todos estes monumentos se escalonam em anfiteatro sobre uma colina
onde está também o nosso amplo \VUgras{liceu} \textsuperscript{(21)}.

%%K t11-c1-07-2 p
Se se desce por uma via um tanto a pique que se junta à Rua Direita,
chega-se ao \VUgras{palácio da justiça} \textsuperscript{(22)}, diante do qual se
estende um agradável \VUgras{jardim público} \textsuperscript{(26)}. Esta
\VUgras{praça ajardinada} não é muito grande, mas põe no meio da cidade um
oásis de verdura e de frescura. Por isso é muito frequentada pelos
habitantes, que gostam de vir sentar-se sob o toldo do \VUgras{café}
\textsuperscript{(25)} para escutar a banda militar que toca às quintas-feiras e
aos domingos no \VUgras{coreto} \textsuperscript{(27)} posto entre os relvados e os
maciços.

%%K t11-c1-08-1 p
8. --- Num desses relvados ergue-se uma \VUgras{estátua} \textsuperscript{(28)} de
bronze, monumento erguido à memória de um dos nossos conterrâneos que
prestou grandes serviços à sua cidade natal.

%%K t11-tit-3 sub
\VUsaut{4.0mm}
\VUpk{10.2pt}{40}{Os transportes.}
\VUsaut{3.4mm}

%%K t11-c1-09-1 p
9. --- A nossa cidade ainda não é iluminada a luz elétrica; só tem
\VUgras{candeeiros a gás} \textsuperscript{(29)}. Mas a \VUgras{imprensa local} faz
campanha para que esse sistema de iluminação se melhore, \VUgras{tal como}
já se aperfeiçoou o \VUgras{sistema de tração} dos elétricos. \VUgras{Os}
velhos carros \VUgras{puxados} por \VUgras{cavalos} foram substituídos por
práticos \VUgras{elétricos} \textsuperscript{(31)}, que levam rapidamente os
viajantes de uma ponta da \VUgras{cidade à outra} por um preço muito
\VUgras{barato}.

%%K t11-c1-10-1 p
10. --- Junto ao palácio da justiça está o \VUgras{banco} \textsuperscript{(32)},
onde se descontam os valores comerciais. Os \VUgras{contínuos do banco}
\textsuperscript{(33)} vão a domicílio fazer as cobranças.

%%K t11-tit-4 sub
\VUsaut{4.0mm}
\VUpk{10.2pt}{40}{Arte e Instrução.}
\VUsaut{3.4mm}

%%K t11-c1-11-1 p
11. --- O belo edifício que há ao lado é o \VUgras{museu}. Contém ao mesmo
tempo a \VUgras{biblioteca} \textsuperscript{(34)} da cidade, belas \VUgras{coleções
de ciências naturais} \textsuperscript{(35)} (museu de história natural) e uma
graciosa \VUgras{galeria de pintura} \textsuperscript{(36)} que constitui uma
esplêndida \VUgras{exposição} permanente.

%%K t11-c1-11-2 p
Está aberto ao público duas vezes por semana, mas os forasteiros podem
visitá-lo em qualquer dia dando uma gorjeta ao \VUgras{guarda}
\textsuperscript{(37)}. Os \VUgras{pintores} \textsuperscript{(38)} vêm ali trabalhar.
Veem-se \VUgras{diante} do seu cavalete, com a \VUgras{paleta} na mão, a
copiar os quadros célebres.

%%K t11-c1-12-1 p
12. --- Na altura, atrás do museu, distingue-se a \VUgras{cúpula}
\textsuperscript{(40)} do \VUgras{hospital} \textsuperscript{(39)} e os edifícios da
\VUgras{escola normal} \textsuperscript{(41)} na qual se formavam os professores
das nossas escolas municipais. Na praça do Teatro há uma \VUgras{estação
dos correios} \textsuperscript{(43)} instalada numa \VUgras{casa particular}
\textsuperscript{(42)}.

%%K t11-tit-5 sub
\VUsaut{5.0mm}
\VUpk{11.4pt}{40}{O Incêndio.}
\VUsaut{3.0mm}

%%K t11-tit-6 sub
\VUpk{10.2pt}{40}{O grito de «Fogo!».}
\VUsaut{3.4mm}

%%K t11-c2-01-1 p
1. --- Uma notícia alarmante espalha-se pela cidade: acaba de declarar-se
um incêndio no teatro! Quando chego à \VUgras{praça} \textsuperscript{(96)}, a
multidão já se apinha no \VUgras{passeio} \textsuperscript{(46)} e na
\VUgras{calçada} \textsuperscript{(47)}. Os \VUgras{empregados dos correios}
\textsuperscript{(44)} e os \VUgras{carteiros} \textsuperscript{(45)} saem da estação. Um
\VUgras{condutor de elétrico} \textsuperscript{(48)} teve de parar o seu carro para
deixar passar uma \VUgras{ambulância} \textsuperscript{(50)} municipal que chega com
um \VUgras{cirurgião} \textsuperscript{(52)} e um \VUgras{enfermeiro}
\textsuperscript{(51)} para socorrer um \VUgras{ferido} \textsuperscript{(53)} que
trouxeram numa \VUgras{maca} \textsuperscript{(54)} até perto do \VUgras{quiosque de
jornais} \textsuperscript{(55)}.

%%K t11-c2-02-1 p
2. --- Uma companhia de infantaria que voltava do exercício parou para
manter a ordem com os \VUgras{guardas municipais} \textsuperscript{(61)} a cavalo.
Os \VUgras{cavaleiros}, cujos cavalos se empinam, fazem-no melhor do que
os \VUgras{soldados} \textsuperscript{(57)} ou os \VUgras{gendarmes}
\textsuperscript{(63)} para fazer recuar os \VUgras{curiosos} \textsuperscript{(56)}. De
resto, estes últimos estão muito ocupados. Um deles indica aos
\VUgras{agentes da polícia} \textsuperscript{(64)} para onde é preciso levar outra
\VUgras{vítima} \textsuperscript{(65)}, um valente bombeiro que caiu de uma escada.

%%K t11-c2-03-1 p
3. --- O \VUgras{oficial} \textsuperscript{(66)} precisa o sítio onde deve
colocar-se a \VUgras{bomba a vapor} \textsuperscript{(67)} que vai chegando.
Enquanto espera essa potente máquina, mandou montar uma \VUgras{bomba de
mão} \textsuperscript{(68)}, cuja longa \VUgras{mangueira} \textsuperscript{(69)} lança
torrentes de água sobre o fogo. Seis bombeiros manejam-na enquanto o seu
companheiro, que segura a \VUgras{agulheta} \textsuperscript{(70)}, dirige o
\VUgras{jato} \textsuperscript{(71)} ao centro do incêndio.

%%K t11-c2-04-1 p
4. --- Do outro lado do teatro encostou-se a grande \VUgras{escada}
\textsuperscript{(72)}. Permite aos bombeiros aproximar-se das chamas que brotam
entre remoinhos de \VUgras{fumo} \textsuperscript{(75)} negro.

%%K t11-tit-7 sub
\VUsaut{5.0mm}
\VUpk{10.2pt}{40}{Atos de valor.}
\VUsaut{3.4mm}

%%K t11-c2-05-1 p
5. --- O \VUgras{incêndio} \textsuperscript{(74)} nasceu no vestíbulo do
\VUgras{teatro} \textsuperscript{(73)}, enquanto se ensaiava a \VUgras{ópera} cuja
primeira \VUgras{representação} devia ter lugar amanhã. Por sorte só havia
uns poucos \VUgras{espetadores} \textsuperscript{(76)} na sala. Os pobres
refugiaram-se na \VUgras{varanda} \textsuperscript{(77)}, aonde valentes
\VUgras{bombeiros} \textsuperscript{(86)} acorrem a socorrê-los com uma escada e
cordas.

%%K t11-c2-06-1 p
6. --- Sob o \VUgras{pórtico} \textsuperscript{(78)}, onde o \VUgras{cartaz}
\textsuperscript{(79)} anuncia a representação, veem-se fugir os \VUgras{músicos}
\textsuperscript{(81)} da \VUgras{orquestra} \textsuperscript{(80)} levando os seus
instrumentos: \VUgras{violino} \textsuperscript{(81)}, \VUgras{flauta}
\textsuperscript{(82)}, \VUgras{violoncelo} \textsuperscript{(83)}, \VUgras{trombone}
\textsuperscript{(84)}, \VUgras{tambor} \textsuperscript{(85)}, etc. Tenta-se salvar parte
do \VUgras{mobiliário} \textsuperscript{(87)}.

%%K t11-c2-07-1 p
7. --- Os \VUgras{atores} \textsuperscript{(89)} e as \VUgras{atrizes}
\textsuperscript{(88)}, \VUgras{cantores} eles e elas, tiveram de fugir com o seu
\VUgras{traje} \textsuperscript{(90)} de teatro. O tenor explica a um
\VUgras{jornalista} \textsuperscript{(92)} como aconteceu. \VUgras{O repórter}
acorreu a correr desde o primeiro momento com as autoridades: o
\VUgras{governador} \textsuperscript{(91)}, o \VUgras{presidente da câmara}
\textsuperscript{(93)}, o \VUgras{comissário de polícia} \textsuperscript{(94)} e um
\VUgras{juiz} \textsuperscript{(95)}, que abrirão um inquérito para determinar as
responsabilidades.
\VUsaut{6.0mm}
\VUfilet{22mm}
